\chapter*{Introduction}
\markright{\textsc{Introduction}}
\addcontentsline{toc}{chapter}{Introduction}
\label{chap:intro}

Environmental sensor networks have become essential infrastructure for understanding natural and built environments.
They operate across a wide range of domains: river gauging stations and catchment areas track water level, discharge,
and quality; urban installations measure air pollution, noise, and microclimate in cities; agricultural plots record
soil moisture, nutrient levels, and crop growth; biodiversity surveys deploy acoustic and camera-trap sensors in
forests and wetlands; and unmanned aerial vehicles carry multispectral or LiDAR payloads over landscapes that are
difficult to instrument permanently. Each of these domains produces continuous streams of heterogeneous data from
devices manufactured by different vendors, communicating over different protocols, and structured according to
different data models.

Managing this diversity at scale, from data acquisition in the field through ingestion, quality control, storage,
and visualization, is a non-trivial engineering problem. Environmental research institutions, national park
authorities, water management agencies, and university departments all face the same challenge: integrating
sensors across geographically distributed sites into a single, standards-compliant data infrastructure.

Czech University of Life Sciences Prague (CZU) is one such institution. As a university focused on life sciences,
agriculture, and environmental research, CZU operates sensor infrastructure spanning river gauging stations,
catchment monitoring networks, and experimental field plots, with instruments recording water level, conductivity,
temperature, turbidity, and other environmental parameters. Prior to this work, no unified platform existed to manage
this infrastructure in a coherent, standards-compliant way.

The UFZ Time.IO solution, developed at the Helmholtz Centre for Environmental Research in Leipzig, was identified as a
candidate foundation. Rather than building a full environmental data management system from scratch, this thesis
investigates whether Time.IO can be adapted to meet the operational requirements of a real institutional deployment,
and implements the extensions needed to close the gap between the upstream system and those requirements.

The platform must also be positioned within the broader context of the eLTER research
infrastructure, for which Time.IO serves as the time-series data acquisition component.
The CZU deployment therefore serves a dual purpose: it addresses a concrete operational
need at CZU, and it provides evidence for or against the suitability of the adapted
platform for wider eLTER adoption.

\medskip\noindent
The goals of this thesis are:
\begin{enumerate}
  \item To analyze the capabilities and limitations of the UFZ Time.IO system with respect to the CZU and eLTER
        deployment requirements.
  \item To design and implement the Water Data Platform (repository \texttt{water-dp}),
        an application layer providing project-oriented sensor management,
        MQTT-driven alert evaluation, geospatial integration, and a multilingual user portal.
  \item To implement a unified deployment layer (repository \texttt{hydro-deploy})
        enabling single-command installation of the complete integrated stack.
  \item To evaluate the resulting platform against the identified requirements and reflect on its suitability for
        broader adoption within the eLTER infrastructure.
\end{enumerate}

\medskip\noindent
The thesis follows a progression from understanding the problem domain to building and validating a solution.
Chapter~\ref{chap:background} first establishes the necessary context: the landscape of environmental sensor data
management, the OGC SensorThings API standard on which the platform relies, the upstream UFZ Time.IO system that
serves as the foundation, the institutional context linking the CZU deployment to the eLTER network through the
Time.IO adoption, and the eLTER research infrastructure whose data integration requirements constrain the design. This background
is essential because the design decisions in later chapters depend directly on the capabilities and constraints of
these existing technologies.

Building on this foundation, Chapter~\ref{chap:requirements} translates the CZU operational needs into concrete
requirements and performs a gap analysis against the baseline Time.IO system. The gaps identified here determine
the scope of the implementation work that follows.

Chapter~\ref{chap:architecture} then presents the architecture designed to close those gaps, explaining how
the two runtime layers, the inherited TSM ingestion infrastructure and the new
data platform application layer, interact and why they are structured as separate but integrated components.
Chapter~\ref{chap:implementation} justifies the major technology selections and describes the
realization of this architecture, including the unified deployment layer.

Finally, Chapter~\ref{chap:testing} evaluates the resulting platform against the requirements defined earlier,
verifying that the identified gaps have been addressed and assessing the platform's readiness for production
deployment.
Chapter~\ref{chap:conclusion} summarizes the contributions, reflects on limitations, and outlines future work.
